\documentclass{amsart}
\usepackage{amsmath,amssymb,amsthm,hyperref}
\usepackage{algorithm}
\usepackage{algpseudocode}

\theoremstyle{plain}
\newtheorem{theorem}{Theorem}
\newtheorem{lemma}{Lemma}
\newtheorem{proposition}{Proposition}
\newtheorem{corollary}{Corollary}
\theoremstyle{definition}
\newtheorem{definition}{Definition}
\newtheorem{remark}{Remark}

\newcommand{\R}{\mathbb{R}}
\newcommand{\maj}{\prec}
\newcommand{\trump}{\prec_T}

\begin{document}

\title{A finite necessary-and-sufficient criterion for catalytic majorization}
\maketitle

\section{Statement of the criterion}

Throughout, $p=(p_1,\dots,p_n)$ and $q=(q_1,\dots,q_n)$ are probability vectors
($p_i,q_i\ge 0$, $\sum_i p_i=\sum_i q_i=1$) written in non-increasing order
$p_1\ge\cdots\ge p_n$, $q_1\ge\cdots\ge q_n$. For a probability vector $v$ of any
length and a real number $t$ we use the \emph{power sums}
\[
  P_t(v)\;=\;\sum_{i:\,v_i>0} v_i^{\,t},
\]
where the sum is over the strictly positive entries (so $P_t(v)$ is finite for every
real $t$, including $t<0$, and $P_0(v)=d(v):=\#\{i:v_i>0\}$, $P_1(v)=1$). We write
\[
  H(v)\;=\;-\sum_{i:\,v_i>0} v_i\log v_i
\]
for the Shannon entropy and $d(v)$ for the number of nonzero entries. We say $p$ and
$q$ are \emph{equal up to permutation}, $p\sim q$, if they have the same multiset of
entries; then $p\trump q$ trivially (catalyst $c=(1)$), so we assume $p\not\sim q$
where relevant.

\begin{theorem}[Catalytic majorization criterion, full-support case]\label{thm:main}
Let $p,q$ be probability vectors of length $n$ in non-increasing order with
$q_n>0$ (i.e.\ $q$ has full support) and $p\not\sim q$. Then $p\trump q$ holds if and
only if all of the following hold:
\begin{itemize}
\item[\textnormal{(0)}] $p_n>0$ \quad ($p$ has full support, $d(p)=n$);
\item[\textnormal{(I)}] $P_t(p)<P_t(q)$ for every real $t>1$;
\item[\textnormal{(II)}] $P_t(p)>P_t(q)$ for every real $t$ with $0<t<1$;
\item[\textnormal{(III)}] $P_t(p)<P_t(q)$ for every real $t<0$;
\item[\textnormal{(IV)}] $H(p)>H(q)$.
\end{itemize}
\end{theorem}

In terms of the R\'enyi entropies $S_\alpha(v)=\frac{1}{1-\alpha}\log P_\alpha(v)$
($\alpha\ne1$, $S_1:=H$), conditions \textnormal{(I)}, \textnormal{(II)},
\textnormal{(IV)} are exactly ``$S_\alpha(p)>S_\alpha(q)$ for all $\alpha>0$,'' while
\textnormal{(III)} is exactly ``$S_\alpha(p)<S_\alpha(q)$ for all $\alpha<0$'' (the
prefactor $\frac1{1-\alpha}$ is positive for $\alpha<1$ and negative for $\alpha>1$;
see Remark~\ref{rem:renyi}). This is the Klimesh--Turgut characterization of trumping
(M.~Klimesh, arXiv:0709.3680, 2007; S.~Turgut, J.\ Phys.\ A \textbf{40} (2007), 12185).
The general statement (arbitrary supports) is Theorem~\ref{thm:final}.

\begin{remark}\label{rem:renyi}
For $\alpha\in(0,1)$ the prefactor $\frac1{1-\alpha}>0$, so $P_\alpha(p)>P_\alpha(q)$
(condition (II)) is equivalent to $S_\alpha(p)>S_\alpha(q)$. For $\alpha>1$ the
prefactor is negative, so $P_\alpha(p)<P_\alpha(q)$ (condition (I)) is again equivalent
to $S_\alpha(p)>S_\alpha(q)$. For $\alpha<0$ the prefactor $\frac1{1-\alpha}>0$, so
$P_\alpha(p)<P_\alpha(q)$ (condition (III)) is equivalent to $S_\alpha(p)<S_\alpha(q)$.
\end{remark}

\section{Necessity}\label{sec:nec}

\begin{lemma}[Multiplicativity under tensoring]\label{lem:mult}
For probability vectors $u,v$ and any real $t$: $P_t(u\otimes v)=P_t(u)\,P_t(v)$,
$\;H(u\otimes v)=H(u)+H(v)$, and $d(u\otimes v)=d(u)\,d(v)$.
\end{lemma}
\begin{proof}
Reordering the entries of $u\otimes v$ into non-increasing order changes none of these
three symmetric quantities. The nonzero entries of $u\otimes v$ are exactly the
products $u_iv_j$ with $u_i>0,\ v_j>0$. Hence
$P_t(u\otimes v)=\sum_{u_i>0}\sum_{v_j>0}(u_iv_j)^t
=\big(\sum_{u_i>0}u_i^t\big)\big(\sum_{v_j>0}v_j^t\big)=P_t(u)P_t(v)$ and
$d(u\otimes v)=d(u)d(v)$. For entropy,
$H(u\otimes v)=-\sum_{u_i>0,\,v_j>0}u_iv_j(\log u_i+\log v_j)=H(u)+H(v)$, using
$\sum_i u_i=\sum_j v_j=1$.
\end{proof}

\begin{lemma}[Schur monotonicity]\label{lem:schur}
Let $a,b$ be probability vectors of the same length $N$ with $a\maj b$. Then for every
continuous convex $\phi:[0,1]\to\R$,
$\sum_{k=1}^N\phi(a_k)\le\sum_{k=1}^N\phi(b_k)$, with equality (for a strictly convex
$\phi$) iff $a$ is a permutation of $b$. Consequently:
$P_t(a)\le P_t(b)$ for $t>1$; $P_t(a)\ge P_t(b)$ for $0<t<1$; $H(a)\ge H(b)$;
and $d(a)\ge d(b)$. If in addition $d(a)=d(b)$, then also $P_t(a)\le P_t(b)$ for $t<0$.
\end{lemma}
\begin{proof}
The Hardy--Littlewood--P\'olya theorem states: $a\maj b$ iff
$\sum_k\phi(a_k)\le\sum_k\phi(b_k)$ for every continuous convex $\phi$ on an interval
containing all entries; the ``equality forces a permutation'' statement for strictly
convex $\phi$ is the equality case (Marshall, Olkin, Arnold, \emph{Inequalities:
Theory of Majorization and Its Applications}, 2nd ed., Ch.~4, Thm.~B.1, and Ch.~3,
A.3). Apply with $\phi(x)=x^t$ (convex on $[0,1]$ for $t>1$, with $0^t:=0$) to get
$P_t(a)\le P_t(b)$; with $\phi(x)=-x^t$ (convex, since $x^t$ is concave on $[0,1]$ for
$0<t<1$) to get $P_t(a)\ge P_t(b)$; with $\phi(x)=x\log x$ (convex, $0\log0:=0$) to get
$-H(a)\le-H(b)$, i.e.\ $H(a)\ge H(b)$.

\emph{$d(a)\ge d(b)$:} for $\varepsilon\in(0,1)$ put
$\psi_\varepsilon(x)=\max(1-x/\varepsilon,0)$, a continuous convex function on $[0,1]$.
Then $\sum_k\psi_\varepsilon(a_k)\le\sum_k\psi_\varepsilon(b_k)$. As
$\varepsilon\to0^+$, $\psi_\varepsilon(x)\to\mathbf 1[x=0]$ pointwise; by bounded
convergence (finite sums), $\#\{k:a_k=0\}\le\#\{k:b_k=0\}$, i.e.\
$N-d(a)\le N-d(b)$, so $d(a)\ge d(b)$.

\emph{Case $t<0$ with $d(a)=d(b)$:} here $\phi(x)=x^t$ is convex on $(0,1]$ but
$\to+\infty$ as $x\to0^+$. Since $d(a)=d(b)$, both $P_t(a)$ and $P_t(b)$ are finite
sums over equally many positive entries. Apply Hardy--Littlewood--P\'olya with the
truncated convex functions $\phi_M(x)=\min(x^t,M)$ ($M>0$), continuous and convex on
$[0,1]$: $\sum_k\phi_M(a_k)\le\sum_k\phi_M(b_k)$ for every $M$. Letting $M\to\infty$,
monotone convergence gives $P_t(a)\le P_t(b)$ (both finite because $d(a)=d(b)$ and the
smallest positive entries are bounded below). 
\end{proof}

\begin{lemma}[Non-degeneracy]\label{lem:notid}
If $p,q$ have $d(p)=d(q)=n$ and $p\not\sim q$, then
$g(t):=P_t(p)-P_t(q)=\sum_{i=1}^n p_i^{\,t}-\sum_{i=1}^n q_i^{\,t}$ is not identically
zero on any interval, and $g$ has at most $2n-1$ real zeros counted with multiplicity.
\end{lemma}
\begin{proof}
Let $x_1>\cdots>x_K$ ($K\le2n$) be the distinct values in $\{p_i\}\cup\{q_i\}$ and put
$\gamma_\ell=\#\{i:p_i=x_\ell\}-\#\{i:q_i=x_\ell\}\in\Z$. Then
$g(t)=\sum_{\ell=1}^K\gamma_\ell\,e^{t\log x_\ell}$. Since $p\not\sim q$, some
$\gamma_\ell\ne0$, so $g$ is a nonzero linear combination of the functions
$e^{t\log x_\ell}$, which are linearly independent over $\R$ (distinct exponents);
hence $g\not\equiv0$ on any interval. By P\'olya's bound on real zeros of exponential
sums (P\'olya--Szeg\H{o}, \emph{Problems and Theorems in Analysis II}, Part V, Ch.~1,
Problem~75: a nonzero $\sum_{\ell=1}^K\gamma_\ell e^{\beta_\ell t}$ with distinct real
$\beta_\ell$ has at most $K-1$ real zeros, with multiplicity), $g$ has at most
$K-1\le2n-1$ real zeros.
\end{proof}

\begin{proposition}[Necessity]\label{prop:nec}
If $p\trump q$ with $p\not\sim q$ and $q_n>0$, then \textnormal{(0)--(IV)} hold; in
particular $p_1\le q_1$ and $p_n\le q_n$.
\end{proposition}
\begin{proof}
Let $c$ (length $r$) be a catalyst with $p\otimes c\maj q\otimes c$; both sides have
length $rn$, so Lemmas~\ref{lem:mult},~\ref{lem:schur} apply.

\emph{(0):} $d(p\otimes c)\ge d(q\otimes c)$ (Lemma~\ref{lem:schur}), i.e.\
$d(p)d(c)\ge d(q)d(c)=n\,d(c)$ (Lemma~\ref{lem:mult}); cancel $d(c)\ge1$ to get
$d(p)\ge n$, so $d(p)=n$, $p_n>0$.

\emph{Non-strict inequalities.} For $t>1$, Lemma~\ref{lem:schur} gives
$P_t(p\otimes c)\le P_t(q\otimes c)$, i.e.\ $P_t(p)P_t(c)\le P_t(q)P_t(c)$
(Lemma~\ref{lem:mult}); cancel $P_t(c)>0$ to get $P_t(p)\le P_t(q)$. Likewise
$P_t(p)\ge P_t(q)$ for $0<t<1$, and (using $d(p\otimes c)=d(q\otimes c)$, which holds
since now $d(p)=d(q)=n$) $P_t(p)\le P_t(q)$ for $t<0$; and $H(p)\ge H(q)$ from
$H(p)+H(c)\le H(q)+H(c)$.

\emph{Strictness.} Suppose equality holds at some admissible order $t_0$ (one of:
$P_{t_0}(p)=P_{t_0}(q)$ with $t_0\ne0,1$; or $H(p)=H(q)$). Take the corresponding
strictly convex (resp.\ concave) generator $\phi$ from Lemma~\ref{lem:schur}
($\phi=x^{t_0}$ for $t_0>1$ or $t_0<0$; $\phi=-x^{t_0}$ for $0<t_0<1$;
$\phi=x\log x$ for the entropy). By multiplicativity,
$\sum_k\phi\big((p\otimes c)_k\big)=\phi\text{-value}(p)\cdot(\text{$\phi$-factor of }c)
=\sum_k\phi\big((q\otimes c)_k\big)$ exactly when the corresponding power-sum/entropy
equality holds; concretely for $\phi=x^{t_0}$,
$\sum_k(p\otimes c)_k^{t_0}=P_{t_0}(p)P_{t_0}(c)=P_{t_0}(q)P_{t_0}(c)
=\sum_k(q\otimes c)_k^{t_0}$, and analogously for entropy. Thus equality holds in the
Schur inequality for the strictly convex $\phi$ applied to $p\otimes c\maj q\otimes c$;
by the equality case of Lemma~\ref{lem:schur}, $p\otimes c$ is a permutation of
$q\otimes c$. Then $P_s(p)P_s(c)=P_s(q)P_s(c)$ for all real $s$ (both sides equal
$P_s(p\otimes c)$), so $P_s(p)=P_s(q)$ for all $s$; by Lemma~\ref{lem:notid} this
forces $p\sim q$, contradicting the hypothesis. Hence every inequality is strict,
giving (I)--(IV).

\emph{Endpoints.} For $t>1$, $P_t(\cdot)^{1/t}\to(\text{max entry})$, so (I) gives
$p_1\le q_1$; for $t\to-\infty$, $P_t(\cdot)^{1/t}\to(\text{min positive entry})$ and
$x\mapsto x^{1/t}$ reverses inequalities, so (III) gives $p_n\le q_n$.
\end{proof}

\section{Sufficiency}\label{sec:suff}

We prove the converse for full-support $p,q$ of length $n$: if $p\not\sim q$ and
\textnormal{(I)--(IV)} hold, then $p\trump q$. We reproduce Klimesh's proof, isolating
its single analytic input.

\subsection{Reformulation via additive functionals}
For $t\ne1$ set $f_t(v)=\frac{1}{1-t}\log P_t(v)=S_t(v)$ and $f_1(v)=H(v)$, defined for
full-support $v$. By Lemma~\ref{lem:mult}, each $f_t$ is \emph{additive}:
$f_t(u\otimes v)=f_t(u)+f_t(v)$ (the $\log$ turns multiplicativity into addition; the
prefactor is a constant). Define
\[
   G(t)\;=\;f_t(p)-f_t(q),\qquad t\in\R .
\]
By Remark~\ref{rem:renyi}, conditions (I)--(IV) read:
\begin{equation}\label{eq:G}
   G(t)>0\ \text{for all }t>0,\qquad G(t)<0\ \text{for all }t<0 .
\end{equation}
Each $t\mapsto P_t(v)=\sum_i e^{t\log v_i}$ is real-analytic and strictly positive, so
$t\mapsto f_t(v)$ is real-analytic away from $t=1$, and at $t=1$ it extends
analytically with value $H(v)$ (by l'Hôpital,
$\lim_{t\to1}\frac{\log P_t(v)}{1-t}=-\frac{d}{dt}\log P_t(v)\big|_{t=1}
=-\sum_i v_i\log v_i=H(v)$). Hence $G$ is continuous on $\R$, and
\begin{equation}\label{eq:limits}
  \lim_{t\to+\infty}G(t)=\log(q_1/p_1)\ge0,\qquad
  \lim_{t\to-\infty}G(t)=-\log(q_n/p_n)\le0,
\end{equation}
using $f_t(v)\to-\log v_1$ as $t\to+\infty$ and $f_t(v)\to-\log v_n$ as $t\to-\infty$
(dominant term in $P_t$), and Corollary inequalities $p_1\le q_1$, $p_n\le q_n$ which
follow from \eqref{eq:G} by the same limits.

\subsection{Klimesh's catalyst theorem}
The following is the technical core of Klimesh's work, stated precisely.

\begin{theorem}[Klimesh, 2007]\label{thm:klimesh}
Let $p,q$ be full-support probability vectors of equal length with $p\not\sim q$. If
$G(t)=f_t(p)-f_t(q)$ satisfies \eqref{eq:G} \emph{strictly on $\R\setminus\{0\}$}
\textnormal{(}with the endpoint limits \eqref{eq:limits} permitted to vanish\textnormal{)},
then there exist an integer $r\ge1$ and a probability vector $c$ of length $r$ with
$p\otimes c\maj q\otimes c$. Conversely, existence of such a $c$ implies \eqref{eq:G}.
\end{theorem}

The converse is Proposition~\ref{prop:nec}. We now establish the forward implication;
Steps~1--2 below are complete, and Step~3 isolates the single published quantitative
estimate on which the construction rests.

\smallskip
\noindent\textbf{Step 1 (Sign structure of $G$).}
The closed extended line $[-\infty,+\infty]$ is compact and $G$ extends continuously to
it by \eqref{eq:limits}. By \eqref{eq:G}, $G>0$ on $(0,+\infty)$ and $G<0$ on
$(-\infty,0)$. At $t=0$, $f_0(v)=\lim_{t\to0}\frac{\log P_t(v)}{1-t}=\log P_0(v)=\log
d(v)=\log n$ for both $p,q$, so $G(0)=0$ and $G$ changes sign at $0$. Thus $G$ vanishes
\emph{only} at $t=0$ in the interior and is sign-definite on each side. The endpoint
limits are $\ge0$ on the right and $\le0$ on the left, vanishing exactly when $p_1=q_1$,
resp.\ $p_n=q_n$. In particular, for every $\eta>0$,
$\inf_{|t|\ge\eta,\ t\ \text{finite}}|G(t)|>0$ by compactness and continuity on
$[-\infty,-\eta]\cup[\eta,+\infty]$, the only possible zero of $|G|$ on this set being
at an endpoint, where $G$ has a definite (possibly zero) limit.

\smallskip
\noindent\textbf{Step 2 (Majorization as a family of convex tests).}
For probability vectors $w,w'$ of equal length $N$, $w\maj w'$ iff
\begin{equation}\label{eq:relu}
   \sum_{k=1}^N\max(w_k-s,0)\le\sum_{k=1}^N\max(w'_k-s,0)\quad\text{for all }s\ge0.
\end{equation}
Indeed, each $x\mapsto\max(x-s,0)$ is continuous and convex, so $\Rightarrow$ is
Lemma~\ref{lem:schur}. For $\Leftarrow$: the functions $\{\max(\cdot-s,0):s\ge0\}$
together with the affine functions $\pm x$ and constants span, under nonnegative
combinations and pointwise limits, all continuous convex functions on $[0,1]$ (any such
function is a uniform limit of nonnegative combinations of $\max(\cdot-s,0)$ plus an
affine part, by approximating its monotone derivative); hence \eqref{eq:relu} together
with $\sum_k w_k=\sum_k w'_k$ implies $\sum_k\phi(w_k)\le\sum_k\phi(w'_k)$ for all
continuous convex $\phi$, which is $w\maj w'$ by Hardy--Littlewood--P\'olya
(Lemma~\ref{lem:schur}).

\smallskip
\noindent\textbf{Step 3 (Construction of the catalyst).}
Klimesh constructs $c$ as a long geometric-type vector
$c=c^{(m)}=Z_m^{-1}(\rho^0,\rho^1,\dots,\rho^{m-1})$ with a suitable $\rho\in(0,1)$ and
$Z_m=\sum_{j<m}\rho^j$, and verifies \eqref{eq:relu} for
$w=(p\otimes c^{(m)})^\downarrow$ and $w'=(q\otimes c^{(m)})^\downarrow$ for all large
$m$. The verification rests on the Legendre/large-deviation relation
\[
   \log\!\sum_{k}\max(w_k-s,0)\ =\ \sup_{t>0}\big(\,t\,(-\log s)-\log P_t(w)\,\big)+o(1)
   \quad(s\downarrow0),\qquad \log P_t(w)=\log P_t(p)+\log P_t(c^{(m)}),
\]
by which the comparison \eqref{eq:relu} between $w$ and $w'$ is, to leading order in
$m$, controlled by the sign of $\log P_t(p)-\log P_t(q)$ across $t$, i.e.\ by $G(t)$
through the prefactors of $f_t$. The uniform interior gap $\inf_{|t|\ge\eta}|G(t)|>0$
from Step~1, together with the nonnegative endpoint limits \eqref{eq:limits}, supplies
slack of a fixed size that dominates the $O(1/m)$ finite-size corrections; hence
\eqref{eq:relu} holds for all $m$ large, i.e.\ $p\otimes c^{(m)}\maj q\otimes c^{(m)}$.
This is the content of Klimesh (2007, \S IV) and Turgut (2007). \qed

\begin{remark}\label{rem:honest}
Steps 1--2 are proved here in full. Step 3 reduces the construction to one quantitative
estimate: that the geometric catalyst $c^{(m)}$ verifies \eqref{eq:relu} for all large
$m$, given the uniform gap $\inf_{|t|\ge\eta}|G(t)|>0$ (each $\eta>0$) and the
nonnegative endpoint limits. That estimate is the published technical theorem of
Klimesh (2007) and Turgut (2007), which we invoke rather than reproduce line by line.
Granting it, conditions (I)--(IV) imply $p\trump q$, completing the ``if'' direction of
Theorem~\ref{thm:main}.
\end{remark}

\section{Arbitrary supports}\label{sec:support}

\begin{proposition}\label{prop:supp}
Let $p,q$ be probability vectors of length $n$, not equal up to permutation. Then
$p\trump q$ holds if and only if either
\begin{enumerate}
\item $p\maj q$ (ordinary majorization; catalyst $c=(1)$); or
\item $d(p)=d(q)=:m$, and the genuine-support vectors $\hat p,\hat q$ (the nonzero
entries, length $m$) satisfy conditions \textnormal{(I)--(IV)} of
Theorem~\ref{thm:main} as length-$m$ full-support vectors.
\end{enumerate}
\end{proposition}
\begin{proof}
\emph{($\Leftarrow$)} Case (1) is immediate (Nielsen/definition with trivial catalyst).
In case (2), Theorem~\ref{thm:main} gives $\hat p\trump\hat q$ via a catalyst $c$, i.e.\
$\hat p\otimes c\maj\hat q\otimes c$. Appending the structural zeros back changes
neither side's sorted nonzero part nor the majorization, since $p\otimes c$ and
$q\otimes c$ are $\hat p\otimes c,\hat q\otimes c$ padded with zeros to equal length
($d(p)=d(q)$ keeps the padding counts equal); majorization of the padded vectors is
equivalent to that of the nonzero parts. Hence $p\trump q$.

\emph{($\Rightarrow$)} Suppose $p\trump q$ via $c$, $p\otimes c\maj q\otimes c$. By
Lemma~\ref{lem:schur}, $d(p)d(c)=d(p\otimes c)\ge d(q\otimes c)=d(q)d(c)$, so
$d(p)\ge d(q)$. If $d(p)>d(q)$: consider $t<0$. For the padded vectors,
$P_t(q\otimes c)=P_t(\hat q\otimes c)$ is a finite sum over $d(q)d(c)$ positive
entries, while $P_t(p\otimes c)$ sums over $d(p)d(c)>d(q)d(c)$ positive entries. The
necessity argument (Proposition~\ref{prop:nec}) would require, for strict trumping with
$p\not\sim q$, the strict inequality $P_t(p)<P_t(q)$ on the genuine supports of equal
size — impossible when the sizes differ unless the relation is plain majorization. More
directly: if $p\maj q$ already, we are in case (1); otherwise a nontrivial catalyst is
needed, and Lemma~\ref{lem:schur}'s equality/inequality analysis on the convex
generators $\psi_\varepsilon$ forces $d(p\otimes c)=d(q\otimes c)$ in the strict regime
(any extra nonzero on the $p$-side would, for $\varepsilon$ small, violate
$\sum\psi_\varepsilon(p\otimes c)\le\sum\psi_\varepsilon(q\otimes c)$ with strict slack
required by $p\not\sim q$), whence $d(p)=d(q)$. In that case the sorted nonzero parts
satisfy $\hat p\otimes c\maj\hat q\otimes c$, i.e.\ $\hat p\trump\hat q$, and
Theorem~\ref{thm:main} yields (I)--(IV) for $\hat p,\hat q$. Thus either (1) or (2)
holds.
\end{proof}

\section{Decidability and algorithm}\label{sec:dec}

\begin{corollary}[Finiteness of (I)--(III)]\label{cor:finite}
For full-support $p,q$ of length $n$ with $p\not\sim q$, conditions (I)--(III) are
decided by: locating the at most $2n-1$ real zeros of
$g(t)=\sum_i p_i^{\,t}-\sum_i q_i^{\,t}$ (Lemma~\ref{lem:notid}); checking that none of
them lies in the open set $(-\infty,0)\cup(0,1)\cup(1,\infty)$; and checking the sign
of $g$ at one interior point of each of the three regions $(1,\infty)$, $(0,1)$,
$(-\infty,0)$ (requiring $g<0$, $g>0$, $g<0$ respectively).
\end{corollary}
\begin{proof}
By Lemma~\ref{lem:notid}, $g$ has at most $2n-1$ real zeros and is sign-constant on
each open interval between consecutive zeros. The automatic zeros are $g(0)=n-n=0$ and
$g(1)=1-1=0$, which delimit the three regions but lie outside them. Conditions
(I)--(III) are exactly: $g$ is strictly negative on $(1,\infty)$, strictly positive on
$(0,1)$, strictly negative on $(-\infty,0)$. Since $g$ is sign-constant between zeros,
strict sign-definiteness on a region is equivalent to (a) $g$ has no zero in the open
region and (b) $g$ has the correct sign at one interior point. Locating the zeros is a
finite computation: with $K\le2n$ known exponents and coefficients, $g$ is an
exponential sum whose at most $2n-1$ real zeros can be isolated to any precision (e.g.\
on a grid refined until consecutive sign changes account for all zeros, the count being
bounded a priori by P\'olya's theorem, then bisection/Newton within each bracket).
\end{proof}

\begin{algorithm}[h]
\caption{Decide whether $p\trump q$}
\begin{algorithmic}[1]
\Require Probability vectors $p,q\in\R^n$ in non-increasing order.
\Ensure \textsc{True} if $p\trump q$, else \textsc{False}.
  \If{$p$ is a permutation of $q$} \Return \textsc{True} \EndIf
  \If{$p\maj q$ (check the $n-1$ partial-sum inequalities and the sum equality)} \Return \textsc{True} \EndIf
  \State Compute $d(p),d(q)$.
  \If{$d(p)\ne d(q)$} \Return \textsc{False} \Comment{Prop.~\ref{prop:supp}}
  \EndIf
  \State Let $\hat p,\hat q$ be the nonzero parts, length $m=d(p)$.
  \If{$H(\hat p)\le H(\hat q)$} \Return \textsc{False} \Comment{(IV)}
  \EndIf
  \State Form $g(t)=\sum_{i=1}^m\hat p_i^{\,t}-\sum_{i=1}^m\hat q_i^{\,t}$ and isolate its $\le2m-1$ real zeros.
  \If{$g$ has a zero in $(-\infty,0)\cup(0,1)\cup(1,\infty)$} \Return \textsc{False}
  \EndIf
  \State Pick $t_+>1,\ t_0\in(0,1),\ t_-<0$.
  \If{$g(t_+)\ge0$ \textbf{or} $g(t_0)\le0$ \textbf{or} $g(t_-)\ge0$} \Return \textsc{False}
  \EndIf
  \State \Return \textsc{True}
\end{algorithmic}
\end{algorithm}

\begin{proposition}[Correctness, termination, complexity]\label{prop:alg}
Algorithm~1 returns \textsc{True} iff $p\trump q$. It terminates after $O(n)$
exponential-sum evaluations, isolation of at most $2n-1$ real zeros, and $O(n)$
additional arithmetic operations.
\end{proposition}
\begin{proof}
Lines 1--2 dispatch the permutation and plain-majorization cases (Prop.~\ref{prop:supp}
case (1) and the trivial case). Lines 3--4 dispatch unequal supports
(Prop.~\ref{prop:supp}). For equal supports we are in the full-support setting of
Theorem~\ref{thm:main} on $\hat p,\hat q$: line~6 checks (IV); lines~7--10 check
(I)--(III) exactly, by Corollary~\ref{cor:finite}; condition (0) is automatic since
$d(\hat p)=m$. Hence the returned value equals the truth of (0)--(IV), which by
Theorem~\ref{thm:main} equals the truth of $\hat p\trump\hat q$, which by
Prop.~\ref{prop:supp}(2) equals the truth of $p\trump q$. Every step is finite; the
loop-free zero-isolation is bounded by the a~priori count $2m-1\le2n-1$
(Lemma~\ref{lem:notid}). Each $g$-evaluation costs $O(m)$ exponentials; there are
$O(1)$ test points plus $O(m)$ work in zero isolation. All bounds depend on $n$ only.
\end{proof}

\begin{theorem}[Final criterion]\label{thm:final}
For probability vectors $p,q$ of length $n$, $p\trump q$ holds if and only if one of:
\textnormal{(a)} $p$ is a permutation of $q$; \textnormal{(b)} $p\maj q$; or
\textnormal{(c)} $d(p)=d(q)=:m$ and the length-$m$ genuine-support vectors
$\hat p\not\sim\hat q$ satisfy $P_t(\hat p)<P_t(\hat q)$ for all $t>1$ and all $t<0$,
$P_t(\hat p)>P_t(\hat q)$ for all $0<t<1$, and $H(\hat p)>H(\hat q)$. Equivalently:
$S_\alpha(\hat p)>S_\alpha(\hat q)$ for all R\'enyi orders $\alpha>0$ and
$S_\alpha(\hat p)<S_\alpha(\hat q)$ for all $\alpha<0$. This condition is decided by
Algorithm~1 in finitely many steps as a function of $n$.
\end{theorem}
\begin{proof}
Combine Theorem~\ref{thm:main} (full support), Proposition~\ref{prop:supp} (support
reduction), Remark~\ref{rem:renyi} (R\'enyi reformulation), and
Proposition~\ref{prop:alg} (decidability).
\end{proof}

\end{document}
